\documentclass[14pt]{extarticle}

\usepackage{preamble}

\begin{document}
    \thispagestyle{empty}
    {
        \setstretch{1.0}
        \centerline{Title head}
    }
    
    \vspace{14pt}

    \begin{center}\bf Title Page\end{center}
    
    \vspace{7pt}

    \begin{center}
        \begin{tabular}{ p{7cm} p{-1.5cm} p{7cm} }
            {\bf Author:} & & Last-name Name Middle-name (russian "ФИО" format)\\
            & & \vspace{10pt}\rule{6cm}{0.019cm}\\ % Signature
            & & «\rule{1cm}{0.019cm}»\ \rule{3cm}{0.019cm}\ 20\_\_ г.\\ % Date
        \end{tabular}
    \end{center}
    
    \vfill
    
    \centerline{Place (e.g. city) year}\vspace{-0,7cm}

    \newpage
    
    \tableofcontents

    \newpage
    
    \phantomsection
    \section*{\normalfont\centerline{\makebox[-2.18cm]{}First Section}}
    \addcontentsline{toc}{section}{\makebox[1.2cm]{}First section} % It sets up an appearance of First Section in the Contents 

    Text text text...

    \newpage

    \section{Second section}
    
    \subsection{Subsection of the second one}
        
    Text text text and citation \cite{wang2024a}...
 
    \begin{figure}[h]
        \centering
        \includegraphics[scale=0.5]{she1-2652401-large}
        \caption{Figure's caption}
    \end{figure}
    
    % Some math
    \begin{equation}
        \label{heateq}
        c\rho\dfrac{\partial T}{\partial t} = \mathrm{div}\lt \lambda \nabla T \rt  + q(\mathbf{r}, t)
    \end{equation}
    
    % These two lines add the reference list title. It is a fallback for \defbibheading, which is in the preamble.sty file.
    %\section*{\hspace{-1.25cm}\normalfont\centerline{СПИСОК ИСПОЛЬЗОВАННЫХ ИСТОЧНИКОВ}}
    %\addcontentsline{toc}{section}{СПИСОК ИСПОЛЬЗОВАННЫХ ИСТОЧНИКОВ}
    
    \newpage
    \printbibliography[heading=bibintoc] % Add the reference list.
\end{document}
\bye
